% TOP_PROBLEM: {"id": 3412, "title": "Slice-ribbon problem", "category_id": 7, "category": {"id": 7, "name": "topology", "display_name": "Topology", "description": "Properties preserved under continuous deformations.", "slug": "topology", "order_index": 7, "created_at": "2026-07-31T15:26:25.671Z"}, "status": "open"}
% FIRST_POSTED: 2026-09-25
% ENTRY_KIND: statement_only
% !TeX program = lualatex
% Definition and sources only; this file is not an LLM solution attempt.
\documentclass[11pt]{article}
\usepackage[margin=25mm]{geometry}
\usepackage{fontspec}
\setmainfont{Latin Modern Roman}
\usepackage{amsmath,amssymb,mathtools}
\usepackage{unicode-math}
\setmathfont{Latin Modern Math}
\usepackage{xurl,hyperref}
\hypersetup{colorlinks=true,urlcolor=blue}
\setcounter{secnumdepth}{0}
\setlength{\parindent}{0pt}
\setlength{\parskip}{5pt}
\setlength{\emergencystretch}{3em}
\begin{document}
\section*{\texorpdfstring{Slice-ribbon problem}{Slice-ribbon problem}}\label{3412}
\subsection{Definitions and mathematical statement}
A knot is a smooth embedding $K:S^1\hookrightarrow S^3$, considered up to ambient isotopy. It is \emph{smoothly slice} if it bounds a smooth properly embedded disk $D^2\hookrightarrow B^4$. A \emph{ribbon singularity} of an immersed disk in $S^3$ is a transverse double arc whose two preimage arcs consist of one interior arc and one properly embedded arc ending on the boundary. A ribbon knot bounds an immersed disk with only these singularities. The assertion is
\[
 K\text{ smoothly slice}\quad\Longrightarrow\quad K\text{ ribbon}.
\]
Equivalently, one asks whether every smoothly slice knot admits a slice disk with no interior local maxima for the radial Morse function, after choosing an appropriate disk. The category is smooth, not merely locally flat topological sliceness.
\par\smallskip\textit{Formulation and concept references:} \href{https://www.openproblemgarden.org/op/slice_ribbon_problem_0}{[S1]}.

\subsection{Short English statement}
Must every knot that bounds a smooth disk in four-dimensional space also bound a disk in three-dimensional space with only ribbon-type self-intersections?
\subsection{Sources}
\begin{itemize}
\item[S1]Fox, R. H. "Some problems in knot theory," in Topology of 3-manifolds and related topics (Proc. The Univ. of Georgia Institute, 1961), pp. 168-176, 1962.
\url{https://www.openproblemgarden.org/op/slice_ribbon_problem_0}.
\textit{Formulation source.}
\end{itemize}
\end{document}
